\chapter{人的灵魂}

\section{一、引言}
人的起源和组成究竟如何？本文从圣经所启示的有关人的起源，来看人性的真理和应用。这是人论的一个重要内容。

\subsection{人是神的创造}
圣经清楚地阐明了人类的起源，就是耶和华神用尘土造人（参创2:7）。人来自尘土，而且人体组成的化学元素在土壤中都可以找到。神是按照自己的形象造人（参创1:27）。神的形象不是物质的。因此，人的本质不在乎组成身体的尘土，乃在乎神吹一口生气，就是活的生命之气，进入人体，由此，人被称之为“有灵的活人”（参创2:7）。圣经明显地启示：人由身体和灵魂两个部分组成，成为一个整体的活人。

\subsection{}
第二，虽然圣经中没有经文直接说：“人由灵魂和身体两部分组成”，但是新旧约中的诸多经文明显地这样教导。因此，在《亚他拿修信经》的第37条中，就根据圣经的教导总结说：身体和灵魂组成一个人。

为什么要相信大公信经呢？因为这些信条都是合乎圣经的真理，可以用来抵抗各种异端的教导。又例如，虽然圣经没有经文直接说“神是三位一体的”，但是圣经的大量经文启示了这一真理。所以，《亚他拿修信经》详细地定义和总结了“神是三位一体”的教义。

正统教义和信经的产生不是盲目的，乃是教会在同各样异端邪教的属灵争战中产生的，阻挡了不合乎圣经的异端或异教思想的大量扩散。

我们要根据圣经的经文，参照大公信经的教义，来认识人的起源和组成。这是本文的主要目的。不合圣经真理的教导，违反大公信经的教义，会导致各种谬误和异端。早期的诺斯底主义异端，有关“魂”的教导不合圣经，到今天仍然滋生出各种异端和邪说，例如，有些假教师传讲所谓的内在医治，宣扬“斩断魂结”的心灵医治。这些都不是圣经的教导，不仅欺骗和伤害信徒，而且严重危害教会的健康成长。

\section{二、人的起源}
经上说：“耶和华神用地上的尘土造人，将生气吹在他鼻孔里，他就成了有灵的活人，名叫亚当。”（创世记2:7）

ESV的英文翻译：The LORD God formed the man of dust from the ground and breathed into his nostrils the breath of life, and the man became a living creature. (Genesis 2:7, ESV)

神用尘土造人的身体，并将生气吹在他里面，就成为有灵的活人。人的身体来自“物质的”尘土，灵魂来自“非物质的”神的“吹气”。这里，“有灵的活人”原文直译为：“人就成为一个活的魂”(KJV的英文翻译: man became a living soul)。ESV的翻译更直接：那人就成为一个活的被造物。(the man became a living creature)虽然人包括物质的身体和非物质的灵魂（生命的气息）两个部分，但二者不能截然分割，乃是一个整体。因为这里的经文，并没有直接说神用地上的泥土造人的身体（body），乃是说造了这个人（man）。

另一方面，活的魂(living soul)亦可指整个活人，就是生命的个体(individual)，就是这个人(man)。例如，经上说：“雅各家来到埃及的共有七十人。”（创46:27下）这里“七十人”原文就是七十个魂。ESV直接翻译为：七十个人。(seventy persons)

因此，圣经中讲“灵魂的救恩”不是指人的某个部分得救赎，乃是指整个人的救赎。有人将人的得救分割为三部份：身体的得救、魂的得救和灵的得救，显然是不合圣经的。

中国俗话说：人活一口气。那一口气没有了，人就没有了。圣经上说，当灵魂离开身体，这个人就死了。从人的创造和起源来看，圣经清楚地教导我们：人是由身体和灵魂两部分组成。

传道书上说：“尘土仍归于地，灵仍归于赐灵的神。”(传12:7)神造人包括两部分：用了尘土造身体，并赐给人有灵魂。在圣经中，灵和魂两个字经常交替和平行使用；汉语中也成为一个词组，是同一个意思。

\section{三、圣经中关于“灵与魂”的教导}
第一，灵与魂没有本质区别。在圣经中，“魂”（psuche/soul）与“灵”(pneuma/spirit)，还包括心思／意念(heart/mind)，都指人的非物质层面，没有本质区别。所以，圣经中的灵与魂可以交替使用、平行使用，是同意词。例如，马利亚说：我心尊主为大，我灵以神我的救主为乐。(路1:46-47，这里“心”原文为“魂”)

又如，圣经中描写死亡时，或说人的魂走了(参创35:18;王上17:21;徒15:26)，或者说人的灵不在了(参诗31:5;路23:46;徒7:59)。灵或魂离开身体，就是死亡。

第二，圣经中，两处经文有灵和魂同时出现的例子。（参帖前5:23;来4:12）但是，权威圣经学者一致认为：这两处经文都不是要区分灵与魂的不同，乃是要特别强调整个人，是重复加强语气。

第一处经文，希伯来书的作者说：神的话像两刃的剑刺入、剖开灵与魂，骨节与骨髓。（参来4:12)意思是说，灵魂本来是一个整体、分不开的，但神的话仍然可以穿透那不可分开和辨明的，且要表达后面经文的意思：连心中分不开的思念和主意，神也能辨明。换句话说，如果灵魂本来就是两部份、是分开的，那也就用不着去刺入和剖开了。

第二处经文，当保罗说：愿你们的灵与魂与身子得蒙保守（参帖前5:23），这里的“灵魂体”同时出现也是表示整个人的加强语气，并没有要将灵与魂分成两部分的意思。况且，保罗一贯的教导是：人是由身体和灵魂两部分组成的。（参罗8:10;林前5:5;7:34;林后7:1;弗2:3;西2:5）因此，神学家约翰慕理坚定地说：有人以为保罗在帖前5:23的经文，要给我们提供一套人性组成部分的定义，是毫无圣经根据的。

第三，有人认为  “灵”的用法多与神有关；而“魂”的用法多与世界有关，所以灵、魂是不同质的，应该分开。他们的显著例子是林前2:14所说的“属血气的人”与“属灵的人”，经文说：“然而属血气的人不领会神圣灵的事，反倒以为愚拙。并且不能知道，因为这些事惟有属灵的人纔能看透。”他们认为，这里的“属血气的人”就是“属魂的人”。

然而，在林前2:14的原文中，“属血气的人”乃是自然的人（natural man），是没有重生悔改的人。而且，保罗在这里论到“属灵人”时，“灵”并非指人性中的那个灵，乃是指着圣灵说的，即重生的人有圣灵内住，是属于圣灵的。保罗说：“如果神的灵住在你们心里，你们就不属肉体，乃属圣灵了。人若没有基督的灵，就不是属基督的。”（罗8:9）所以，保罗所说的属血气的人就是没有圣灵内住的非基督徒，而属灵人就是指有圣灵内住的基督徒。

正是因为所有重生的基督徒都是属灵人，有圣灵内住，我们不是用人智慧所指教的言语，乃是用圣灵所指教的言语，将属灵的话解释属灵的事。（参林前2:13）所以，保罗才总结说：“属灵的人能看透万事，却没有一人能看透了他。”（林前2:15）因此，这里的“属灵人”乃是说“属圣灵的人”。这句经文根本就不支持灵与魂应该分开的假说。

约翰慕理说：当“属灵”一词用在事物上时，意味着来源于圣灵或被圣灵所创造；当用在人身上时，意味着被圣灵所内住、引导、和掌管。

虽然圣经没有证据支持灵魂分开的假说。不过，我们不必假设魂和灵在文法上总是完全相同、可以互换的。魂和灵指示的同一件事可以从不同的角度来看。一种文法情况下，“灵”是恰当的用法，另一种情况下，“魂”则更恰当。例如，死亡被说成是放弃灵，却也被说成是搁下魂（生命）。不同的情况用不同的词，却表达同一个意思。这是经文字义的丰富。圣经的依据如下。

一则，圣经中，描写死人复活时，既可以说魂返回入身体（参王上17:22），也可以说灵回来了。（参路8:55）灵魂的用法是很广的。

二则，圣经说：神是个“灵”（约4:24）。但是，“魂”也可以用来指耶和华神。例如，以赛亚书42:1节：“看哪，我（耶和华）的仆人，我所扶持，所拣选，心里所喜悦的”（耶和华的“心”原文为魂）。还有其它类似经文。（参耶9:9;来10:38;摩6:8）

这里特别提醒，有人教导说，将“灵”定义为高尚的、好的，而“魂”是低级的、恶的。所以，要操练治死魂、提升灵。这种教导完全没有圣经根据，而且有亵渎的嫌疑。

《亚他拿修信经》是早期大公教会四大信经之一，曾为马丁路德所钟爱。该信经第三十七条中明确认定圣经的真理为：人是由灵魂和身体两部分组成的。大量经文证明大公信经的教导是合乎圣经的。一千五百多年来，正统的基督教会都认同这一信条。

\section{四、灵与魂不可分}
首先是解经的问题。如果灵与魂同时出现，就表示二者有本质的不同，那么如何处理其它类似经文呢？其一，圣经上说：“你要尽心、尽性、尽力爱耶和华你的神。”（申6:5），心、性、力，再加上体，那就变成四元了。其二，当主耶稣说：“你要尽心、尽性、尽意、尽力，爱主你的神。”（可12:30）这里，心、性、意、力，再加上体，这难道是说人由五个部分组成？显然不是！解经不能这样望文生义的，应该考察圣经的原文字义、上下文理和全本圣经的系统教导，并且要遵守解释圣经的总原则。

其次，世俗哲学的影响。中国教会有一种传统将“魂”解释为支配人的心思意念和情感等，这是源于古希腊哲学的思想，圣经上并没有这种“魂”的教导。神是个灵，难道神没有思想和情感吗？经上说：耶和华不仅发怒，也施恩、好怜悯、行公义。神是智慧的源头。（参出4:14;赛30:18;诗19:7,111:10）“魂”乃是指人的非物质层面，或整个人的生命，或者活物、位格等。魂在圣经中的用法非常广泛。

再次，圣经中有许多灵和魂的不可分的情况，举例如下。

第一例。马利亚说：“我心（原文：魂）尊主为大，我灵以神我的救主为乐。”（路1:46-47）圣经明确教导：马利亚的魂、灵都是尊主为大的，以神为乐的。如果马利亚的魂不尊主为大，马利亚的灵也不会以神为乐的。因为灵和魂是不分的，均是指同一个人的。

第二例。有人以为：动物有魂无灵，人有魂又有灵的。由此教导说：人必须去掉魂，才能减少动物的成分即兽性。这种说法蛮有理性的，但不合圣经的教导。因为人失掉魂，就意味身体死亡、生命的终结。(参创35:18;王上17:21;徒15:26)

其实，这种“动物有魂无灵”的圣经依据是：“谁知道人的灵是往上升，兽的魂是下入地呢。”（传道书3:21）然而，希伯来原文中兽的“魂”与人的“灵”都是同一个字“灵”(spirit)。英文圣经都正确地译为“兽的灵”。这也是和合本的翻译不足之处。

另外，古希腊哲学以为属物质的身体是邪恶的，而人的灵是善的。这种教导也是不合圣经的真理。因为圣经强调自亚当堕落后，人是全然败坏，包括人的身心灵之全然污秽。（参诗53:3;罗1:24,3:12）神救赎人，也是全人救赎，不仅拯救灵魂，也要在末日使身体复活，即拯救身体。经上说：“神救赎我的灵魂免入深坑；我的生命也必见光。”（约伯记33:28）

保罗所教导的“属灵人”乃是指属圣灵的，而不是指人性中的那个“灵”。（参罗8:9）属灵人就是指有圣灵内住的基督徒。（参罗8:16）所以，属灵人不是操练出来的，而是神在基督里所赐给我们白白的恩典。基督徒靠圣灵的能力得救，也是靠圣灵的能力过成圣的生活。基督徒是世界的光，要效法基督，过荣神益人的生活。

灵与魂分开的错误教导，为某些推动新纪元运动的人所采用。例如，有些搞“内在医治”的“牧师”说，“斩断魂结”是必须的心灵医疗过程。他们将“魂结”定义为一个人与另一个人紧密而长久的连系，到死才结束。这种连系包括“魂”的三个领域：思想、情感和意志。造成不好的“魂结”的原因是多方面的，例如，性关系，家庭暴力，亲戚朋友不健康的关系，上瘾行为等等。他们认为这些“魂结”是犯罪的原因，所以要斩断魂结，就可以除去罪的捆绑。

显然，圣经没有“魂结”和“斩断魂结”的教导。这种教导将人的罪归于环境和人际关系。这是一种世俗心理学或异教巫术的东西。圣经教导我们：人犯罪是因为我们是罪人、有罪性，而非因有坏的“魂结”。“斩断魂结”的操练不但不能帮助人，反而使人陷入更深的罪恶。基督徒要认罪悔改、过成圣的生活，讨神的喜悦。


\section{五、结语}

神用地上的泥土造人，又吹一口生气，将灵魂或生命赐给人，人就成为有理性有情感的活物。这是人的起源，也是灵魂的起源。身体和灵魂乃是一个整体，不可绝然分开。神是按照祂自己的形象造人的。人的堕落，乃是整个人的堕落，是身心灵的堕落。神要拯救人，恢复伊甸园的创造。祂不仅要救人的灵魂，而且在基督再来时，死去信徒的身体也要复活得赎。人由身体和灵魂两个部分组成。灵和魂是同质不可分的，用法上也是可互换的。身体代表人的物质层面，灵魂代表人的非物质层面。这是神的话语启示，也是圣经的真理。

基督徒除了熟读圣经以外，学习和明白大公教会的信条和正统的神学教义也是必须的。这样就可以预防和抵制各种异端邪说。因为大公信条和正统的神学教义是来自圣经的真理，是两千年来教会在与异端邪教作斗争的过程中产生的。末后的时代，人们厌烦纯正的道理，喜欢听信假师父的谎言，教会更不应该忽略纯真的教义，要用真道来责备、警戒和劝勉人。（参提后4:2-5）

圣经上的话是不可让人随私意解说的，强解经文的人，是自取沉沦。(参彼后1:20;3:16)传道人要坚守圣经的真理，按正意分解神的道，教导纯正的教义。（参提后1:13;多1:9,2:1）愿人都高举基督，学习正统的神学教义，回归圣经的真理。

 





















